\section{System V AMD64 ABI}

\subsection{本章要解决的问题}

前面几章已经能看懂一段函数内部的汇编：整数指令、地址计算、条件跳转、循环、\instruction{call} 和 \instruction{ret}。但只知道这些还不够。真实程序不是一个函数孤立运行，而是由大量函数、目标文件、静态库、动态库、运行时和操作系统共同组成。

问题来了：一个 C++ 文件里编译出来的函数，为什么能被另一个 C 文件调用？一个手写汇编函数，为什么能返回给 C++ 调用者？动态库里的 \code{printf}，为什么知道你的参数放在哪里？异常展开、调试器栈回溯、profiler 采样，为什么能沿着调用链走？

答案是 \term{ABI}{Application Binary Interface}。ABI 是二进制层面的契约。它规定：

\begin{itemize}
  \item 函数参数放在哪些寄存器或栈位置。
  \item 返回值放在哪里。
  \item 哪些寄存器由调用者保存，哪些由被调用者保存。
  \item 栈在调用点如何对齐。
  \item 结构体、浮点、向量、变参函数如何传递。
  \item 函数入口、函数尾声、调试信息和异常展开如何约定。
\end{itemize}

本章讲 Linux、macOS 等类 Unix x86-64 平台常用的 System V AMD64 ABI。Windows x64 使用另一套 ABI，参数寄存器、shadow space、保存规则都不同。本书后续默认讨论 System V AMD64，除非明确说明。

\begin{keyidea}
ISA 告诉 CPU 一条指令怎么执行；ABI 告诉不同二进制代码如何互相调用。读函数级汇编时，如果不懂 ABI，就无法可靠判断参数、返回值、栈槽、寄存器保存和跨语言边界。
\end{keyidea}

\subsection{ABI 为什么必须存在}

考虑两个源文件：

\begin{lstlisting}[language=C++]
// add.cpp
extern "C" int add2(int a, int b)
{
    return a + b;
}
\end{lstlisting}

\begin{lstlisting}[language=C++]
// main.cpp
#include <cstdio>

extern "C" int add2(int, int);

int main()
{
    std::printf("%d\n", add2(10, 32));
}
\end{lstlisting}

它们可以分开编译：

\begin{lstlisting}[language=bash]
clang++ -std=c++20 -O3 -c add.cpp -o add.o
clang++ -std=c++20 -O3 -c main.cpp -o main.o
clang++ add.o main.o -o app
\end{lstlisting}

\code{main.o} 编译时并不知道 \code{add2} 的机器码是什么，但它知道怎样调用 \code{add2}：

\begin{lstlisting}
edi = 10
esi = 32
call add2
read return value from eax
\end{lstlisting}

\code{add.o} 编译时也不知道谁会调用它，但它知道入口时参数在哪里：

\begin{lstlisting}
on entry:
    edi = first int argument
    esi = second int argument

on return:
    eax = int return value
\end{lstlisting}

这就是 ABI 的意义。调用者和被调用者不需要共享源代码，不需要由同一个编译器同时优化，只要遵守同一套二进制约定，就可以链接和运行。

\begin{mentalmodel}
函数调用不是“跳到另一个 C++ 函数”。机器层面更像一次协议交接：
\begin{lstlisting}
caller:
    arrange arguments
    preserve caller-owned live values if needed
    align stack
    call callee
    read return value

callee:
    assume arguments are in ABI-defined locations
    preserve callee-saved registers it modifies
    restore stack
    ret
\end{lstlisting}
\end{mentalmodel}

\subsection{先掌握最常见的整数参数}

System V AMD64 ABI 中，前 6 个整数或指针类参数通常放在这些寄存器里：

\begin{center}
\begin{tabular}{p{0.16\linewidth}p{0.25\linewidth}p{0.39\linewidth}}
\toprule
参数序号 & 64 位寄存器 & 常见 32 位整数视图 \\
\midrule
第 1 个 & \register{rdi} & \register{edi} \\
第 2 个 & \register{rsi} & \register{esi} \\
第 3 个 & \register{rdx} & \register{edx} \\
第 4 个 & \register{rcx} & \register{ecx} \\
第 5 个 & \register{r8} & \register{r8d} \\
第 6 个 & \register{r9} & \register{r9d} \\
\bottomrule
\end{tabular}
\end{center}

整数返回值通常放在 \register{rax}；如果返回 \code{int}，有效结果在 \register{eax}。

C++：

\begin{lstlisting}[language=C++]
extern "C" int sum6(int a, int b, int c, int d, int e, int f)
{
    return a + b + c + d + e + f;
}
\end{lstlisting}

可能汇编：

\begin{lstlisting}
sum6:
    lea eax, [rdi + rsi]
    add eax, edx
    add eax, ecx
    add eax, r8d
    add eax, r9d
    ret
\end{lstlisting}

入口状态：

\begin{lstlisting}
edi = a
esi = b
edx = c
ecx = d
r8d = e
r9d = f
\end{lstlisting}

返回约定：

\begin{lstlisting}
eax = a + b + c + d + e + f
ret
\end{lstlisting}

\begin{warningbox}
\register{rcx} 在 System V ABI 中是第 4 个整数参数；在 Windows x64 ABI 中，\register{rcx} 是第 1 个整数参数。看资料或反汇编时必须先确认目标 ABI。
\end{warningbox}

\subsection{超过 6 个整数参数：栈上传参}

第 7 个及以后的整数或指针类参数通常通过栈传递。看例子：

\begin{lstlisting}[language=C++]
extern "C" long sum8(long a, long b, long c, long d,
                     long e, long f, long g, long h)
{
    return a + b + c + d + e + f + g + h;
}
\end{lstlisting}

一种可能汇编：

\begin{lstlisting}
sum8:
    lea rax, [rdi + rsi]
    add rax, rdx
    add rax, rcx
    add rax, r8
    add rax, r9
    add rax, qword ptr [rsp + 8]
    add rax, qword ptr [rsp + 16]
    ret
\end{lstlisting}

函数刚进入时，\register{rsp} 指向返回地址：

\begin{lstlisting}
on entry to sum8:

    [rsp + 0]   return address
    [rsp + 8]   7th argument: g
    [rsp + 16]  8th argument: h

    rdi = a
    rsi = b
    rdx = c
    rcx = d
    r8  = e
    r9  = f
\end{lstlisting}

为什么第 7 个参数是 \code{[rsp + 8]} 而不是 \code{[rsp]}？因为 \instruction{call} 已经把返回地址压到了栈顶。被调用函数入口处，栈顶首先是返回地址，栈上传递的第一个参数在返回地址上方 8 字节。

\begin{deepdive}
如果被调用函数一进来就 \instruction{push rbp} 或 \instruction{sub rsp, ...}，参数相对 \register{rsp} 的偏移会改变。因此读汇编时必须知道当前 \register{rsp} 已经移动了多少。使用帧指针时，编译器常用 \register{rbp} 作为稳定基准。
\end{deepdive}

\subsection{浮点和向量参数}

浮点参数走 SSE/AVX 寄存器通道。常见规则：

\begin{itemize}
  \item \code{float}、\code{double} 参数通常放在 \register{xmm0} 到 \register{xmm7}。
  \item 浮点返回值通常放在 \register{xmm0}。
  \item 整数参数寄存器序列和浮点参数寄存器序列分别计数。
\end{itemize}

C++：

\begin{lstlisting}[language=C++]
extern "C" double mix_args(int a, double b, int c, float d)
{
    return double(a + c) + b + double(d);
}
\end{lstlisting}

入口位置通常是：

\begin{lstlisting}
edi  = a
xmm0 = b
esi  = c
xmm1 = d
\end{lstlisting}

注意 \code{c} 是第 2 个整数类参数，所以在 \register{esi}，不是 \register{edx}。因为 \code{double b} 消耗的是浮点寄存器 \register{xmm0}，不消耗整数参数寄存器。

可能汇编片段：

\begin{lstlisting}
mix_args:
    add edi, esi
    cvtsi2sd xmm2, edi
    addsd xmm2, xmm0
    cvtss2sd xmm1, xmm1
    addsd xmm1, xmm2
    movapd xmm0, xmm1
    ret
\end{lstlisting}

返回值在 \register{xmm0}。这里的 \instruction{cvtsi2sd} 把整数转换成 double，\instruction{cvtss2sd} 把 float 扩展成 double。

\begin{warningbox}
\register{xmm0} 既可能是第 1 个浮点参数，也可能是浮点返回值。函数入口和函数返回是两个不同时间点。读汇编时要区分“入口时 \register{xmm0} 是参数”和“返回时 \register{xmm0} 是结果”。
\end{warningbox}

\subsection{返回值位置}

常见返回值约定：

\begin{center}
\begin{tabular}{p{0.28\linewidth}p{0.28\linewidth}p{0.30\linewidth}}
\toprule
返回类型 & 返回位置 & 说明 \\
\midrule
\code{int}、\code{unsigned} & \register{eax} & 写 32 位会清零 \register{rax} 高 32 位 \\
\code{long}、指针 & \register{rax} & 64 位整数或地址 \\
\code{double}、\code{float} & \register{xmm0} & 标量浮点返回 \\
两个整数机器字 & \register{rax}、\register{rdx} & 常见小结构体返回形式 \\
大结构体 & 隐藏 sret 指针 & 调用者提供返回对象地址 \\
\bottomrule
\end{tabular}
\end{center}

小结构体例子：

\begin{lstlisting}[language=C++]
struct Pair {
    long x;
    long y;
};

extern "C" Pair make_pair(long a, long b)
{
    return Pair{a + 1, b + 2};
}
\end{lstlisting}

可能汇编：

\begin{lstlisting}
make_pair:
    lea rax, [rdi + 1]
    lea rdx, [rsi + 2]
    ret
\end{lstlisting}

返回时：

\begin{lstlisting}
rax = result.x
rdx = result.y
\end{lstlisting}

调用者按 ABI 知道这两个寄存器共同组成返回对象。

\subsection{大结构体返回：隐藏 sret 指针}

大结构体通常不能只靠 \register{rax} 和 \register{rdx} 返回。调用者会先分配好返回对象空间，然后把地址作为隐藏参数传给被调用者。

C++：

\begin{lstlisting}[language=C++]
struct Big {
    long a;
    long b;
    long c;
};

extern "C" Big make_big(long x)
{
    return Big{x, x + 1, x + 2};
}
\end{lstlisting}

逻辑上像这样：

\begin{lstlisting}[language=C++]
extern "C" Big make_big(long x);
\end{lstlisting}

机器层面更像：

\begin{lstlisting}[language=C++]
extern "C" void make_big_hidden(Big* out, long x);
\end{lstlisting}

可能汇编：

\begin{lstlisting}
make_big:
    mov rax, rdi
    mov qword ptr [rdi], rsi
    lea rcx, [rsi + 1]
    mov qword ptr [rdi + 8], rcx
    add rsi, 2
    mov qword ptr [rdi + 16], rsi
    ret
\end{lstlisting}

入口状态：

\begin{lstlisting}
rdi = hidden pointer to caller-allocated result object
rsi = x
\end{lstlisting}

返回状态：

\begin{lstlisting}
memory at rdi:
    [rdi + 0]  = x
    [rdi + 8]  = x + 1
    [rdi + 16] = x + 2

rax = rdi
\end{lstlisting}

这解释了一个重要现象：源代码里看起来只有一个参数 \code{x}，但汇编入口时 \register{rdi} 可能不是 \code{x}，而是隐藏返回对象指针。真正的 \code{x} 变成了 \register{rsi}。

\begin{keyidea}
看到函数第一个参数“不对劲”时，不要马上怀疑编译器。先检查返回类型是否是大结构体，是否存在隐藏 sret 指针。
\end{keyidea}

\subsection{结构体参数的简化分类规则}

System V AMD64 ABI 对聚合类型有详细分类算法。完整规则涉及 \code{INTEGER}、\code{SSE}、\code{SSEUP}、\code{X87}、\code{COMPLEX_X87}、\code{MEMORY} 等类别，还会把对象按 8 字节分块。本书先给出足以读大多数性能代码的简化流程：

\begin{enumerate}
  \item 把结构体按 8 字节为单位切成一个或多个 eightbyte。
  \item 如果对象太大、字段未对齐、或 C++ 类型有非平凡拷贝/析构语义，通常走内存或隐藏引用。
  \item 如果不超过两个 eightbyte，且每块能分类成 INTEGER 或 SSE，就可能通过寄存器传递。
  \item INTEGER 类通常走通用寄存器。
  \item SSE 类通常走 \register{xmm} 寄存器。
  \item 一旦需要的寄存器不够，相关参数可能改走栈。
\end{enumerate}

例子 1：

\begin{lstlisting}[language=C++]
struct TwoLong {
    long a;
    long b;
};

extern "C" long use_two_long(TwoLong x)
{
    return x.a + x.b;
}
\end{lstlisting}

两个 8 字节整数块，可能入口为：

\begin{lstlisting}
rdi = x.a
rsi = x.b
\end{lstlisting}

例子 2：

\begin{lstlisting}[language=C++]
struct TwoDouble {
    double a;
    double b;
};

extern "C" double use_two_double(TwoDouble x)
{
    return x.a + x.b;
}
\end{lstlisting}

两个 8 字节 SSE 块，可能入口为：

\begin{lstlisting}
xmm0 = x.a
xmm1 = x.b
\end{lstlisting}

例子 3：

\begin{lstlisting}[language=C++]
struct Mixed {
    int a;
    double b;
};

extern "C" double use_mixed(Mixed x)
{
    return double(x.a) + x.b;
}
\end{lstlisting}

常见入口可能是：

\begin{lstlisting}
edi  = x.a
xmm0 = x.b
\end{lstlisting}

这说明结构体并不总是“整体放到内存”。小结构体可能被拆分到多个寄存器里。

\begin{warningbox}
结构体 ABI 分类是容易出错的地方。字段顺序、padding、对齐、浮点/整数混合、C++ 非平凡类型都会改变传参方式。遇到公开二进制接口时，不要凭感觉改结构体布局。
\end{warningbox}

\subsection{caller-saved 和 callee-saved}

寄存器数量有限。函数调用前后，哪些寄存器可以被破坏？哪些必须保持？ABI 用 caller-saved 和 callee-saved 分工。

\begin{center}
\begin{tabular}{p{0.18\linewidth}p{0.34\linewidth}>{\raggedright\arraybackslash}p{0.34\linewidth}}
\toprule
类别 & 寄存器 & 含义 \\
\midrule
caller-saved & \register{rax}、\register{rcx}、\register{rdx}、\register{rsi}、\register{rdi}、\register{r8}--\register{r11} & 调用者若还需要这些值，调用前自己保存 \\
callee-saved & \register{rbx}、\register{rbp}、\register{r12}--\register{r15} & 被调用者如果修改，返回前必须恢复 \\
特殊 & \register{rsp} & 返回前必须恢复到入口约定状态 \\
浮点/向量 & \register{xmm0}--\register{xmm15} & System V 下通常 caller-saved \\
\bottomrule
\end{tabular}
\end{center}

为什么这样分？如果所有寄存器都由调用者保存，每次调用都很重；如果所有寄存器都由被调用者保存，每个函数入口都很重。ABI 把一部分寄存器定义为 caller-saved，适合短生命周期临时值；把另一部分定义为 callee-saved，适合跨调用仍要保留的值。

\subsubsection{callee-saved 示例}

C++：

\begin{lstlisting}[language=C++]
extern "C" long g(long);

extern "C" long keep_across_call(long x)
{
    long y = x * 3;
    long z = g(x);
    return y + z;
}
\end{lstlisting}

\code{y} 必须跨过对 \code{g} 的调用继续存在。编译器可能把 \code{y} 放在 callee-saved 寄存器 \register{rbx}：

\begin{lstlisting}
keep_across_call:
    push rbx
    lea  rbx, [rdi + rdi*2]
    call g
    add  rax, rbx
    pop  rbx
    ret
\end{lstlisting}

解释：

\begin{lstlisting}
push rbx
    save caller's original rbx

lea rbx, [rdi + rdi*2]
    rbx = x * 3

call g
    g may clobber caller-saved registers
    g must preserve rbx if it uses rbx

add rax, rbx
    rax = g(x) + x * 3

pop rbx
    restore caller's original rbx

ret
\end{lstlisting}

如果 \code{keep_across_call} 修改了 \register{rbx} 却没有恢复，调用它的上层函数可能崩溃或产生错误结果。

\begin{deepdive}
callee-saved 不是说这个寄存器“不能用”。它的意思是：可以用，但你要把调用者看到的旧值还回去。常见保存方式是 \instruction{push} 和 \instruction{pop}，也可以保存到栈帧中的固定偏移。
\end{deepdive}

\subsection{栈对齐规则}

System V AMD64 ABI 要求调用点满足 16 字节栈对齐。更具体地说：

\begin{lstlisting}
before call:
    rsp is 16-byte aligned

call pushes 8-byte return address

on callee entry:
    rsp + 8 is 16-byte aligned
    rsp itself is 8 mod 16
\end{lstlisting}

一个函数入口常见状态：

\begin{lstlisting}
rsp points to return address
rsp mod 16 = 8
\end{lstlisting}

如果函数要再调用别的函数，它必须在自己的 \instruction{call} 前把 \register{rsp} 调整到 16 字节对齐。

例子：

\begin{lstlisting}
caller:
    sub rsp, 8
    call callee
    add rsp, 8
    ret
\end{lstlisting}

假设 \code{caller} 入口时 \code{rsp mod 16 = 8}：

\begin{lstlisting}
entry:
    rsp mod 16 = 8

sub rsp, 8:
    rsp mod 16 = 0

call callee:
    pushes return address
    callee entry rsp mod 16 = 8

add rsp, 8:
    restore caller's stack pointer
\end{lstlisting}

为什么要对齐？一方面 ABI 需要统一约定；另一方面 SIMD load/store、局部对象对齐、调用约定和编译器生成代码都依赖这个规则。一些指令如 \instruction{movaps} 对内存地址有对齐要求；即使现代 CPU 对很多未对齐访问能处理，错误的 ABI 栈对齐仍然可能导致崩溃或性能下降。

\begin{warningbox}
手写汇编调用 C/C++ 函数时，最常见错误之一就是忘记在 \instruction{call} 前对齐 \register{rsp}。这种错误可能在简单测试里不暴露，换一个编译器版本、库函数或 CPU 状态就崩。
\end{warningbox}

\subsection{栈帧、帧指针和局部变量}

带帧指针的函数常见形态：

\begin{lstlisting}
push rbp
mov  rbp, rsp
sub  rsp, 32
...
add  rsp, 32
pop  rbp
ret
\end{lstlisting}

或等价尾声：

\begin{lstlisting}
leave
ret
\end{lstlisting}

进入函数后栈可能长这样：

\begin{lstlisting}
higher addresses

[rbp + 16]  stack argument 2 if present
[rbp + 8]   return address
[rbp + 0]   saved old rbp
[rbp - 8]   local/spill slot
[rbp - 16]  local/spill slot
[rbp - 24]  local/spill slot

lower addresses
\end{lstlisting}

\register{rbp} 的价值是稳定。函数里 \register{rsp} 可能因为 \instruction{push}、\instruction{sub rsp, ...}、调用准备而变化，但 \register{rbp} 通常保持不变，方便调试器、低优化级代码和人工阅读。

优化级较高时，编译器常省略帧指针：

\begin{lstlisting}[language=bash]
clang++ -std=c++20 -O3 -fomit-frame-pointer ...
\end{lstlisting}

性能分析时，很多工程会保留帧指针以改善采样栈回溯：

\begin{lstlisting}[language=bash]
clang++ -std=c++20 -O3 -fno-omit-frame-pointer ...
\end{lstlisting}

是否保留帧指针是性能、调试、profiling 之间的工程权衡。现代生产服务常为了 profiler 可用性保留帧指针。

\subsection{red zone}

System V AMD64 ABI 定义了 \term{red zone}{红区}：用户态函数入口时，\register{rsp} 下方 128 字节区域可以被叶子函数临时使用，而不必移动 \register{rsp}。

位置：

\begin{lstlisting}
[rsp - 128] ... [rsp - 1]   red zone
[rsp + 0]                  return address
\end{lstlisting}

叶子函数是不调用其他函数的函数。例子：

\begin{lstlisting}[language=C++]
extern "C" long leaf(long x)
{
    long tmp[4];
    tmp[0] = x + 1;
    tmp[1] = x + 2;
    tmp[2] = x + 3;
    tmp[3] = x + 4;
    return tmp[0] + tmp[1] + tmp[2] + tmp[3];
}
\end{lstlisting}

如果优化器不能完全消除数组，叶子函数可能使用 \register{rsp} 下方空间，而不执行 \instruction{sub rsp, ...}。

red zone 的限制：

\begin{itemize}
  \item 只适合不调用其他函数的叶子函数。
  \item 一旦执行 \instruction{call}，新的返回地址会写到更低地址，可能覆盖调用者 red zone 中的内容。
  \item 内核、裸机、中断上下文通常不能依赖 red zone，常用 \code{-mno-red-zone}。
  \item 手写汇编使用 red zone 时必须确认目标环境允许。
\end{itemize}

\begin{lstlisting}[language=bash]
clang++ -std=c++20 -O3 -mno-red-zone ...
\end{lstlisting}

\subsection{变参函数和 \code{va_list}}

变参函数比普通函数复杂，因为被调用者不知道静态参数列表之后还有多少参数、类型是什么。

典型调用：

\begin{lstlisting}[language=C++]
#include <cstdio>

extern "C" void print_values(int x, double y)
{
    std::printf("x=%d y=%f\n", x, y);
}
\end{lstlisting}

调用 \code{printf} 时：

\begin{lstlisting}
rdi  = pointer to format string
esi  = x
xmm0 = y
al   = number of vector registers used for floating arguments
call printf
\end{lstlisting}

System V ABI 要求调用变参函数时，\register{al} 携带用于传递浮点/向量参数的向量寄存器数量上界。你经常在调用 \code{printf} 前看到：

\begin{lstlisting}
mov eax, 1
call printf
\end{lstlisting}

如果没有浮点参数，可能看到：

\begin{lstlisting}
xor eax, eax
call printf
\end{lstlisting}

变参函数内部使用 \code{va_start} 时，编译器会维护一个 \code{va_list}，里面通常包含：

\begin{lstlisting}
gp_offset
fp_offset
overflow_arg_area
reg_save_area
\end{lstlisting}

直觉上：

\begin{itemize}
  \item 前几个普通参数可能来自寄存器保存区。
  \item 超出寄存器容量的参数来自栈上的 overflow 区域。
  \item 浮点参数和整数参数有不同 offset。
\end{itemize}

\begin{deepdive}
手写汇编调用变参函数时，不能只把参数放到寄存器就结束。尤其是调用带浮点参数的 \code{printf}，还要正确设置 \register{al}。这个细节非常适合实验，因为错误代码有时能运行，有时输出错，有时在不同 libc 或优化级别下表现不同。
\end{deepdive}

\subsection{C++ 名字修饰和 \code{extern "C"}}

ABI 不只包含寄存器和栈，也包含符号名约定。C++ 支持函数重载、命名空间、类成员函数，因此编译器会做 \term{name mangling}{名字修饰}。

C++：

\begin{lstlisting}[language=C++]
int add(int, int);
double add(double, double);
\end{lstlisting}

两个函数源代码名都叫 \code{add}，目标文件里的符号名必须不同。编译器可能生成类似：

\begin{lstlisting}
_Z3addii
_Z3adddd
\end{lstlisting}

如果希望 C++ 函数使用 C ABI 符号名，可以写：

\begin{lstlisting}[language=C++]
extern "C" int add(int, int);
\end{lstlisting}

这样符号名通常就是 \code{add}。本书实验里大量使用 \code{extern "C"}，不是因为 C++ 不好，而是为了让汇编符号名稳定、易读、易链接。

\subsection{手写汇编函数的最低正确性清单}

如果你要写一个被 C++ 调用的 \code{.S} 函数，至少检查：

\begin{enumerate}
  \item 符号名和 C++ 声明是否一致。
  \item 参数寄存器是否按 System V ABI 读取。
  \item 返回值是否写到正确寄存器。
  \item 修改 callee-saved 寄存器前是否保存，返回前是否恢复。
  \item 返回前 \register{rsp} 是否恢复。
  \item 调用其他函数前 \register{rsp} 是否 16 字节对齐。
  \item 调用变参函数前 \register{al} 是否正确。
  \item 是否误用 red zone。
  \item 是否需要 \code{.type}、\code{.size} 和 \code{.note.GNU-stack}。
\end{enumerate}

一个最小汇编函数：

\begin{lstlisting}
.intel_syntax noprefix
.text
.globl asm_add2
.type asm_add2, @function
asm_add2:
    lea eax, [rdi + rsi]
    ret
.size asm_add2, .-asm_add2
.section .note.GNU-stack,"",@progbits
\end{lstlisting}

对应 C++ 声明：

\begin{lstlisting}[language=C++]
extern "C" int asm_add2(int, int);
\end{lstlisting}

\subsection{案例：C++ 调汇编}

创建两个文件。

\filepath{labs/ch13_system_v_abi/asm_add.S}：

\begin{lstlisting}
.intel_syntax noprefix
.text
.globl asm_add6
.type asm_add6, @function
asm_add6:
    lea eax, [rdi + rsi]
    add eax, edx
    add eax, ecx
    add eax, r8d
    add eax, r9d
    ret
.size asm_add6, .-asm_add6
.section .note.GNU-stack,"",@progbits
\end{lstlisting}

\filepath{labs/ch13_system_v_abi/main.cpp}：

\begin{lstlisting}[language=C++]
#include <cstdio>

extern "C" int asm_add6(int, int, int, int, int, int);

int main()
{
    int r = asm_add6(1, 2, 3, 4, 5, 6);
    std::printf("%d\n", r);
}
\end{lstlisting}

构建：

\begin{lstlisting}[language=bash]
clang++ -std=c++20 -O2 -c main.cpp -o main.o
clang++ -c asm_add.S -o asm_add.o
clang++ main.o asm_add.o -o abi_demo
./abi_demo
\end{lstlisting}

预期输出：

\begin{lstlisting}
21
\end{lstlisting}

反汇编：

\begin{lstlisting}[language=bash]
objdump -d -Mintel abi_demo > abi_demo.objdump.txt
\end{lstlisting}

你应该在 \code{main} 的调用点看到参数被放入 ABI 寄存器，在 \code{asm_add6} 中看到这些寄存器被读取。

\subsection{案例：汇编调 C++}

C++：

\begin{lstlisting}[language=C++]
#include <cstdio>

extern "C" int cpp_square(int x)
{
    return x * x;
}

extern "C" int asm_call_square(int);

int main()
{
    std::printf("%d\n", asm_call_square(9));
}
\end{lstlisting}

汇编：

\begin{lstlisting}
.intel_syntax noprefix
.text
.globl asm_call_square
.type asm_call_square, @function
asm_call_square:
    sub rsp, 8
    call cpp_square
    add eax, 1
    add rsp, 8
    ret
.size asm_call_square, .-asm_call_square
.section .note.GNU-stack,"",@progbits
\end{lstlisting}

为什么要 \code{sub rsp, 8}？因为 \code{asm_call_square} 入口时 \code{rsp mod 16 = 8}。在它执行 \instruction{call} \code{cpp_square} 之前，必须让 \code{rsp mod 16 = 0}。

执行路径：

\begin{lstlisting}
main:
    edi = 9
    call asm_call_square

asm_call_square:
    align stack
    call cpp_square
    eax = cpp_square(9) + 1
    restore stack
    ret
\end{lstlisting}

返回值：

\begin{lstlisting}
cpp_square(9) = 81
asm_call_square(9) = 82
\end{lstlisting}

\subsection{实验 13.1：参数寄存器和栈参数地图}

\begin{labbox}
创建 \filepath{labs/ch13_system_v_abi/args.cpp}：

\begin{lstlisting}[language=C++]
#include <cstdint>

extern "C" long sum8(long a, long b, long c, long d,
                     long e, long f, long g, long h)
{
    return a + b + c + d + e + f + g + h;
}

extern "C" double mix(int a, double b, int c, float d)
{
    return double(a + c) + b + double(d);
}

struct TwoLong {
    long a;
    long b;
};

extern "C" long use_two_long(TwoLong x)
{
    return x.a * 10 + x.b;
}
\end{lstlisting}

生成：

\begin{lstlisting}[language=bash]
mkdir -p labs/ch13_system_v_abi/build
clang++ -std=c++20 -O0 -S -masm=intel args.cpp \
  -o labs/ch13_system_v_abi/build/args_O0.s
clang++ -std=c++20 -O3 -S -masm=intel args.cpp \
  -o labs/ch13_system_v_abi/build/args_O3.s
clang++ -std=c++20 -O3 -c args.cpp \
  -o labs/ch13_system_v_abi/build/args.o
objdump -d -Mintel labs/ch13_system_v_abi/build/args.o \
  > labs/ch13_system_v_abi/build/args.objdump.txt
\end{lstlisting}

交付物：

\begin{enumerate}
  \item 对 \code{sum8} 列出 8 个参数分别在哪里。
  \item 解释为什么第 7、第 8 个参数从栈读取。
  \item 对 \code{mix} 列出整数参数和浮点参数的寄存器序列。
  \item 对 \code{use_two_long} 判断结构体是否拆到寄存器。
  \item 对比 \code{-O0} 和 \code{-O3} 栈帧差异。
  \item 从 \code{objdump} 记录至少 8 条机器码字节。
\end{enumerate}
\end{labbox}

\subsection{实验 13.2：破坏 callee-saved 寄存器}

\begin{labbox}
创建 \filepath{labs/ch13_system_v_abi/callee_saved.S}：

\begin{lstlisting}
.intel_syntax noprefix
.text

.globl bad_clobber_rbx
.type bad_clobber_rbx, @function
bad_clobber_rbx:
    mov rbx, 0x1234567812345678
    ret
.size bad_clobber_rbx, .-bad_clobber_rbx

.globl good_clobber_rbx
.type good_clobber_rbx, @function
good_clobber_rbx:
    push rbx
    mov rbx, 0x1234567812345678
    pop rbx
    ret
.size good_clobber_rbx, .-good_clobber_rbx

.globl check_bad
.type check_bad, @function
check_bad:
    push rbx
    mov rbx, 0x1122334455667788
    call bad_clobber_rbx
    mov rax, rbx
    pop rbx
    ret
.size check_bad, .-check_bad

.globl check_good
.type check_good, @function
check_good:
    push rbx
    mov rbx, 0x1122334455667788
    call good_clobber_rbx
    mov rax, rbx
    pop rbx
    ret
.size check_good, .-check_good

.section .note.GNU-stack,"",@progbits
\end{lstlisting}

创建 \filepath{labs/ch13_system_v_abi/callee_saved.cpp}：

\begin{lstlisting}[language=C++]
#include <cstdint>
#include <cstdio>

extern "C" std::uint64_t check_bad();
extern "C" std::uint64_t check_good();

int main()
{
    std::printf("bad  = 0x%llx\n",
        (unsigned long long)check_bad());
    std::printf("good = 0x%llx\n",
        (unsigned long long)check_good());
}
\end{lstlisting}

构建：

\begin{lstlisting}[language=bash]
clang++ -std=c++20 -O2 -c callee_saved.cpp -o callee_saved.o
clang++ -c callee_saved.S -o callee_saved_asm.o
clang++ callee_saved.o callee_saved_asm.o -o callee_saved_demo
./callee_saved_demo
objdump -d -Mintel callee_saved_demo > callee_saved_demo.objdump.txt
\end{lstlisting}

交付物：

\begin{enumerate}
  \item 解释 \code{bad_clobber_rbx} 违反了哪条 ABI 规则。
  \item 解释 \code{good_clobber_rbx} 如何保存和恢复 \register{rbx}。
  \item 说明 \code{check_bad} 为什么返回被破坏后的值。
  \item 在 GDB 中单步观察 \register{rbx} 变化。
  \item 把这个实验推广到 \register{r12}，写一个新的坏例子和修复例子。
\end{enumerate}
\end{labbox}

\subsection{实验 13.3：栈对齐和 \instruction{movaps}}

\begin{labbox}
创建 \filepath{labs/ch13_system_v_abi/stack_align.S}：

\begin{lstlisting}
.intel_syntax noprefix
.text

.globl bad_movaps_store
.type bad_movaps_store, @function
bad_movaps_store:
    sub rsp, 16
    xorps xmm0, xmm0
    movaps xmmword ptr [rsp], xmm0
    add rsp, 16
    ret
.size bad_movaps_store, .-bad_movaps_store

.globl good_movaps_store
.type good_movaps_store, @function
good_movaps_store:
    sub rsp, 24
    xorps xmm0, xmm0
    movaps xmmword ptr [rsp], xmm0
    add rsp, 24
    ret
.size good_movaps_store, .-good_movaps_store

.section .note.GNU-stack,"",@progbits
\end{lstlisting}

创建 \filepath{labs/ch13_system_v_abi/stack_align.cpp}：

\begin{lstlisting}[language=C++]
#include <cstdio>

extern "C" void bad_movaps_store();
extern "C" void good_movaps_store();

int main()
{
    std::puts("calling good");
    good_movaps_store();
    std::puts("calling bad");
    bad_movaps_store();
    std::puts("done");
}
\end{lstlisting}

构建并运行：

\begin{lstlisting}[language=bash]
clang++ -std=c++20 -O2 -c stack_align.cpp -o stack_align.o
clang++ -c stack_align.S -o stack_align_asm.o
clang++ stack_align.o stack_align_asm.o -o stack_align_demo
./stack_align_demo
\end{lstlisting}

交付物：

\begin{enumerate}
  \item 计算 \code{bad_movaps_store} 中 \instruction{sub rsp, 16} 后的栈对齐状态。
  \item 计算 \code{good_movaps_store} 中 \instruction{sub rsp, 24} 后的栈对齐状态。
  \item 解释 \instruction{movaps} 为什么要求目标地址对齐。
  \item 如果程序崩溃，用 GDB 或 shell 记录信号和崩溃位置。
  \item 把 \instruction{movaps} 改成 \instruction{movups}，观察行为是否变化，并解释为什么不能把这当成“不需要 ABI 对齐”的证据。
\end{enumerate}
\end{labbox}

\subsection{实验 13.4：大结构体返回和隐藏指针}

\begin{labbox}
创建 \filepath{labs/ch13_system_v_abi/sret.cpp}：

\begin{lstlisting}[language=C++]
#include <cstdint>

struct Small {
    std::uint64_t a;
    std::uint64_t b;
};

struct Big {
    std::uint64_t a;
    std::uint64_t b;
    std::uint64_t c;
};

extern "C" Small make_small(std::uint64_t x)
{
    return Small{x, x + 1};
}

extern "C" Big make_big(std::uint64_t x)
{
    return Big{x, x + 1, x + 2};
}
\end{lstlisting}

生成：

\begin{lstlisting}[language=bash]
clang++ -std=c++20 -O0 -S -masm=intel sret.cpp -o sret_O0.s
clang++ -std=c++20 -O3 -S -masm=intel sret.cpp -o sret_O3.s
clang++ -std=c++20 -O3 -c sret.cpp -o sret.o
objdump -d -Mintel sret.o > sret.objdump.txt
\end{lstlisting}

交付物：

\begin{enumerate}
  \item 说明 \code{Small} 如何返回。
  \item 说明 \code{Big} 如何返回。
  \item 找出 \code{make_big} 中隐藏 sret 指针的位置。
  \item 解释为什么源代码参数 \code{x} 在汇编中可能不是第一个参数寄存器。
  \item 画出调用者为 \code{Big} 分配返回对象并传入地址的过程。
\end{enumerate}
\end{labbox}

\subsection{实验 13.5：变参函数和 \register{al}}

\begin{labbox}
创建 \filepath{labs/ch13_system_v_abi/varargs.cpp}：

\begin{lstlisting}[language=C++]
#include <cstdio>

extern "C" void print_int(int x)
{
    std::printf("x=%d\n", x);
}

extern "C" void print_double(double x)
{
    std::printf("x=%f\n", x);
}

extern "C" void print_mix(int x, double y)
{
    std::printf("x=%d y=%f\n", x, y);
}
\end{lstlisting}

生成：

\begin{lstlisting}[language=bash]
clang++ -std=c++20 -O2 -S -masm=intel varargs.cpp -o varargs.s
clang++ -std=c++20 -O2 -c varargs.cpp -o varargs.o
objdump -d -Mintel varargs.o > varargs.objdump.txt
\end{lstlisting}

交付物：

\begin{enumerate}
  \item 找出每次调用 \code{printf} 前 \register{eax} 或 \register{al} 如何设置。
  \item 解释 \code{print_int} 为什么通常设置为 0。
  \item 解释 \code{print_double} 和 \code{print_mix} 为什么通常设置为 1。
  \item 写一个手写汇编函数调用 \code{printf} 打印 double，先故意错误设置 \register{al}，再修复。
  \item 记录错误版本和修复版本在你的系统上的表现。
\end{enumerate}
\end{labbox}

\subsection{本章作业}

\begin{exercisebox}
\subsubsection*{作业 1：ABI 参数表}

对下面函数签名，给出 System V AMD64 下每个参数的常见传递位置：

\begin{lstlisting}[language=C++]
extern "C" long f1(long a, int b, void* p, long d,
                  long e, long f, long g);

extern "C" double f2(int a, double b, int c,
                    float d, double e);

struct P { long a; long b; };
extern "C" long f3(P p, long x);
\end{lstlisting}

要求：

\begin{enumerate}
  \item 区分通用寄存器和 \register{xmm} 寄存器。
  \item 标出哪些参数会走栈。
  \item 对结构体参数说明你的分类依据。
  \item 说明返回值位置。
\end{enumerate}

\subsubsection*{作业 2：栈布局推导}

给定函数入口汇编：

\begin{lstlisting}
push rbp
mov  rbp, rsp
push rbx
sub  rsp, 24
\end{lstlisting}

要求画出执行完这些指令后的栈布局，至少包含：

\begin{itemize}
  \item 返回地址。
  \item saved \register{rbp}。
  \item saved \register{rbx}。
  \item 24 字节局部区域。
  \item \register{rbp} 和 \register{rsp} 分别指向哪里。
\end{itemize}

\subsubsection*{作业 3：修复错误汇编}

下面函数被 C++ 调用：

\begin{lstlisting}
.intel_syntax noprefix
.globl broken
broken:
    mov rbx, rdi
    call helper
    add rax, rbx
    ret
\end{lstlisting}

要求：

\begin{enumerate}
  \item 找出所有 ABI 风险。
  \item 修复 callee-saved 寄存器问题。
  \item 修复调用前栈对齐问题。
  \item 给出修复后的完整汇编。
  \item 解释每条新增指令为什么必要。
\end{enumerate}

\subsubsection*{作业 4：sret 逆向}

给定汇编：

\begin{lstlisting}
make_obj:
    mov rax, rdi
    mov qword ptr [rdi], rsi
    mov qword ptr [rdi + 8], rdx
    mov qword ptr [rdi + 16], rcx
    ret
\end{lstlisting}

要求：

\begin{enumerate}
  \item 判断是否存在隐藏返回对象指针。
  \item 推测源代码返回类型大小。
  \item 推测显式参数可能有哪些。
  \item 解释为什么返回值还会写 \register{rax}。
\end{enumerate}

\subsubsection*{作业 5：跨语言 ABI 小项目}

写一个小项目，包含：

\begin{itemize}
  \item 一个 C++ 文件声明并调用 3 个手写汇编函数。
  \item 汇编函数 1：整数参数求和，至少 8 个参数。
  \item 汇编函数 2：调用一个 C++ helper，并正确保存 callee-saved 寄存器。
  \item 汇编函数 3：返回一个两字段结构体。
  \item CMake 或 Makefile 构建脚本。
  \item 单元测试或简单断言。
  \item \code{objdump} 报告，解释每个函数的 ABI 行为。
\end{itemize}

\subsubsection*{评分标准}

本章作业按下面标准评价：

\begin{enumerate}
  \item 参数位置是否和 ABI 一致。
  \item 是否能区分寄存器参数和栈参数。
  \item 是否能正确处理 caller-saved 与 callee-saved。
  \item 手写汇编是否保持栈对齐。
  \item 是否能识别隐藏 sret 指针。
  \item 是否能用 GDB、\code{objdump} 和实际运行结果证明结论。
  \item 是否明确区分 System V AMD64 和其他 ABI。
\end{enumerate}
\end{exercisebox}

\subsection{本章小结}

ABI 是从“单个函数内部汇编”走向“真实二进制程序”的关键桥梁。本章你应该掌握：

\begin{lstlisting}
integer/pointer args:
    rdi, rsi, rdx, rcx, r8, r9, then stack

floating args:
    xmm0 - xmm7

return:
    rax / rdx for many integer-like results
    xmm0 for scalar floating results
    hidden sret pointer for large aggregate results

register ownership:
    caller-saved: caller protects live values
    callee-saved: callee restores what it modifies

stack:
    on entry, rsp points to return address
    before call, rsp must be 16-byte aligned
    leaf functions may use red zone in allowed environments
\end{lstlisting}

从高性能计算和 AI 算子开发角度，ABI 决定了 kernel 函数边界的成本和正确性：参数如何进入寄存器，哪些寄存器可以自由使用，调用 helper 前要保存什么，返回 tensor metadata 或小结构体会不会产生隐藏内存写。下一章会把这些规则用于 C++ 和汇编互操作，进一步进入可测试、可链接、可维护的手写汇编工程。
